\section{Coherent polar induction}\label{sec:polar}

\subsection{Spaces, markers, and bad schemes}

Fix a field $k$, a two-dimensional $k$-space $E$, and
\[
  S_n=\PP(\Gamma^nE).
\]
For $\lambda\in E^\vee$, define divided-power contraction by
\[
 \langle\iota_\lambda f,g\rangle=\langle f,\lambda g\rangle.
\]
Then
\begin{equation}\label{eq:lift}
 g\in W_{\iota_\lambda f}\quad\Longleftrightarrow\quad
 \lambda g\in W_f.
\end{equation}
The lower witness lifts squarefreely only when $\lambda\nmid g$.
Thus the factor selected for contraction must remain as a marked
forbidden root.

\begin{definition}[Persistent and modular loci]\label{def:carriers}
The \emph{persistent locus} is the ordinary rank-two catalecticant
locus propagated by contraction.  The \emph{modular locus} is the
additional contraction-kernel locus that can appear in positive
characteristic.  A \emph{collision divisor} records parameters at
which the marked factor becomes a repeated root.
\end{definition}

\begin{definition}[Contraction and marker spaces]\label{def:polar-spaces}
Let $C_f:E^\vee\to\Gamma^{n-1}E$ be
$\lambda\mapsto\iota_\lambda f$, and put
$U_f=\PP(E^\vee)\setminus\PP(\ker C_f)$.  The first-polar rational map
is
\[
 \Phi_f:U_f\longrightarrow S_{n-1},\qquad
 [\lambda]\longmapsto[\iota_\lambda f]
.
\]
Its image is a point or a line.  For $j\geq1$, let
\[
\operatorname{Conf}_j(\PP(E^\vee))
=\{(\lambda_1,\ldots,\lambda_j):\lambda_i\ne\lambda_\ell
  \text{ for }i\ne\ell\}.
\]
The ordered marker space $U_{f,j}$ is the open subset on which every
successive contraction is nonzero.  The universal $j$-step polar flag
is the map
\[
(\lambda_1,\ldots,\lambda_j)\longmapsto
\bigl([\iota_{\lambda_1}f],\ldots,
[\iota_{\lambda_j}\cdots\iota_{\lambda_1}f]\bigr).
\]
\end{definition}

In an affine coordinate $r$ the first map is
\[
\Phi_f(r)=(a_1-ra_0,\ldots,a_n-ra_{n-1}),\qquad
\Phi_f(\infty)=(a_0,\ldots,a_{n-1}).
\]

\begin{example}[The complete R5-to-R6 propagation at $q=29$]
\label{ex:q29-r6}
Work over $\F_{29}$ at redundancy six and take
\[
 f=(6:1:0:0:5:1).
\]
The Hankel kernel is the net of quartics
\[
 W_f=\ker
 \begin{pmatrix}6&1&0&0&5\\1&0&0&5&1\end{pmatrix}
 =
 \left\langle
 (0,0,1,0,0),(24,1,0,1,0),(28,1,0,0,1)
 \right\rangle .
\]
The displayed basis has no common factor, so the persistent test does
not decide the syndrome.  Choose the first marker $r=0$.  Then
\[
 \iota_0f=(1:0:0:5:1).
\]
This is now a redundancy-five syndrome, and its full cubic pencil is
\[
 W_{\iota_0f}
 =\left\langle(0,1,0,0),(25,0,1,24)\right\rangle .
\]
The member with pencil parameter \([9:28]\) is
\[
 h(t)=4+9t-t^2+5t^3\in W_{\iota_0f},
\]
because the two Hankel products are $4+5\cdot5=0$ and
$5(-1)+5=0$ in $\F_{29}$.  Direct factorization gives
\[
 h(t)=5(t-2)(t-6)(t-27).
\]
None of these roots is the marked root $0$.  Equation~\eqref{eq:lift}
therefore lifts $h$ to
\[
 t\,h(t)=(0,4,9,28,5)\in W_f,
\]
a squarefree quartic with roots $0,2,6,27$.  The Hankel criterion now
places $f$ in the span of the corresponding four normal-rational-curve
columns, so $f$ is shallow.  This one calculation exhibits every
interface in the R5-to-R6 propagation: the upper Hankel net, the
ordinary common-factor test, contraction to the complete R5 cubic
pencil, selection on its splitting cover, marker avoidance, and
squarefree lifting.  In the general proof the R5 curve argument
supplies such a pencil parameter after deleting the marked fibers; the
displayed parameter \([9:28]\) is one complete instance of that step.
\end{example}

\begin{definition}[One-step lower package]\label{def:lower-package}
Fix \(f\in S_n(\F_q)\) with injective contraction map, so that
\(\Phi_f:\PP(E^\vee)\to S_{n-1}\) is defined everywhere.  A one-step
lower package abstracts the R6 argument: a lower bad locus, the R5
ordered-root curve with its deletions, and the finite set of markers
where that construction is unavailable.  It consists of:
\begin{enumerate}[label=\textup{(\roman*)}]
\item a closed bad scheme \(B_{n-1}\subset S_{n-1}\) and a finite
  stratification \(S_{n-1}\setminus B_{n-1}=\coprod_\alpha X_\alpha\);
\item for every
  \(\lambda\in\PP(E^\vee)(\F_q)\) with
  \(b=\iota_\lambda f\in X_\alpha(\F_q)\), a specified proper
  geometrically integral curve \(Y_{\alpha,b}\), defined over
  \(\F_q\), with genus bound \(g_\alpha\) and point bound
  \[
       \#Y_{\alpha,b}(\F_q)\geq
       q+\kappa_\alpha-2g_\alpha\sqrt q;
  \]
  together with a zero-dimensional closed deletion scheme
  \(D_{\alpha,b,\lambda}\subset Y_{\alpha,b}\), of degree at most
  \(\delta_\alpha\), containing the singular, branch, diagonal,
  marker, and collision points, such that every
  rational point outside \(D_{\alpha,b,\lambda}\) determines an
  ordered rational split squarefree \(g\in W_b\) with
  \(\lambda\nmid g\);
\item a finite parameter scheme \(Z_f\subset\PP(E^\vee)\), of degree
  at most \(d\), containing
  \(\Phi_f^{-1}(B_{n-1})\) and every parameter at which the preceding
  construction is unavailable, including parameters where a fixed
  factor of the lower system equals a retained marker.
\end{enumerate}
The index \(\alpha\) includes any constant-field or Frobenius twist;
the assertion in (iii) is the required descent statement.  All degrees
are geometric degrees on the displayed curve or parameter line.
\end{definition}

The construction is summarized in Figure~\ref{fig:polar-flow}.  The
upper branch is the level-specific component calculation; the lower
branch is the general transverse argument.
\begin{figure}[H]
\centering
\[
\begin{array}{ccccc}
[f]\in S_n&\dashrightarrow&\Phi_f(U_f)\subset S_{n-1}
 &\longrightarrow&\text{lower ordered splitting cover}\\
&&\begin{array}{c}
\searrow\ \text{contained}\\[-2pt]
\nearrow\ \text{transverse}
\end{array}
&&\begin{array}{l}
\text{persistent/Lucas component},\\
\text{delete branch, diagonal, and markers}
\end{array}\\
&&&&\downarrow\ \text{lift by \eqref{eq:lift}}\\
&&&&\text{split squarefree witness in }W_f .
\end{array}
\]
\caption{Coherent polar induction.  The contained branch is a
level-specific component theorem; only the transverse branch is
controlled by the general point-count argument.}
\label{fig:polar-flow}
\end{figure}

\begin{theorem}[Polar-flag construction]\label{thm:polar-construction}
The maps of Definition~\ref{def:polar-spaces} commute with field
extension and the natural $\operatorname{GL}(E)$ action.  They include
the infinity chart displayed above.  If
$g\in W_{\iota_{\lambda_j}\cdots\iota_{\lambda_1}f}$, then
$\lambda_1\cdots\lambda_jg\in W_f$; the lift is squarefree exactly
when $g$ is squarefree and avoids every marker $\lambda_i$.
\end{theorem}

\begin{proof}
Contraction is defined by the perfect divided-power/symmetric-power
pairing, so it commutes with base change and change of basis.  Repeated
application of \eqref{eq:lift} gives the membership assertion.  Since
the $\lambda_i$ are distinct on the configuration space, their product
is squarefree; the product with $g$ is squarefree precisely when $g$
is squarefree and none of the $\lambda_i$ divides it.  The coordinate
formula at infinity follows by evaluating $C_f$ on the second basis
covector.
\end{proof}

\begin{lemma}[Three readings of the collision divisor]
\label{lem:marker-collision}
Assume \(W_f\) has trivial gcd, and let \(\lambda\in\PP(E^\vee)\).
The following are equivalent:
\begin{enumerate}[label=\textup{(\roman*)}]
\item every member of \(W_f\) divisible by \(\lambda\) is divisible
  by \(\lambda^2\);
\item \(\lambda\) is a fixed factor of \(W_{\iota_\lambda f}\);
\item \(\lambda\) lies on the collision divisor of the linear series
  \(W_f\).
\end{enumerate}
\end{lemma}

\begin{proof}
Equation~\eqref{eq:lift} gives the equality of subspaces
\[
 \lambda W_{\iota_\lambda f}
 =W_f\cap\lambda\Sym^{n-2}E^\vee.
\]
Thus (i) and (ii) are the same statement after dividing by
\(\lambda\).  Since \(W_f\) is base-point-free, (i) says exactly that
the morphism defined by \(W_f\) has zero differential at \(\lambda\),
which is (iii).
\end{proof}

\begin{lemma}[Uniform collision alternatives]
\label{lem:uniform-collision}
Let \(j\geq5\), and let a base-point-free
\(g^{\,j-4}_{\,j-2}\) on \(\PP^1\) define an inseparable morphism in
characteristic \(p\).  Then
\[
 p\mid(j-2),\qquad p(j-4)\leq j-2.
\]
The only possible inseparable cases are \((j,p)=(5,3)\) and
\((6,2)\).  In the separable case the collision divisor has degree at
most \(2j-6\); in particular this holds unconditionally for \(j\geq7\).
\end{lemma}

\begin{proof}
An inseparable morphism factors through absolute Frobenius.  Hence
\(p\mid(j-2)\), and its \(j-3\) sections lie in the
\((j-2)/p+1\)-dimensional space of \(p\)-th powers.  This gives the
inequality.  Since \((j-2)/(j-4)<2\) for \(j\geq7\), no prime remains;
the two smaller cases follow directly.  In the separable case some
value--derivative minor is nonzero; as a binary Wronskian it has degree
\(2(j-2)-2=2j-6\) and contains the collision scheme.
\end{proof}

\begin{lemma}[Linear modular pullbacks]
\label{lem:linear-modular-pullback}
Let \(\mathcal N\subset S_{j-2}\) be a projective linear modular
nucleus.  Then \(\Phi_f^{-1}(\mathcal N)\) is empty, one point, or the
whole parameter line.  In the transverse case its degree is at most
one.
\end{lemma}

\begin{proof}
The polar map has linear coordinates in the parameter, so the
pullback ideal is generated by linear forms on \(\PP^1\).
\end{proof}

\begin{definition}[One-step contained-component assertion]
\label{def:contained}
For a specified bad scheme \(B_{n-1}\), collision divisor, and explicit
closed list \(\mathcal C_n\subset S_n\), the assertion
\(\mathrm{CC}(n,1;\mathcal C_n)\) says that every \(f\) for which the
contraction map is noninjective and hence supplies no polar line,
\(\Phi_f^*B_{n-1}\) is the whole parameter line and hence supplies no
transverse marker, or the collision divisor is the whole line and
hence every lift repeats its marked root, belongs to \(\mathcal C_n\).
When the list is clear we write
\(\mathrm{CC}(n,1)\).
\end{definition}

\begin{definition}[Stable-component assertion]\label{def:stable-component}
The assertion \(\mathrm{SC}(j)\) says that the recursively pointed
contained bad scheme at redundancy \(j\) has no component outside
\(\mathcal P_j\cup\mathcal M_j\).  It is a geometric component
statement and contains no field-size threshold.
\end{definition}

The indices follow two conventions.  The assertion
\(\mathrm{CC}(n,1;\mathcal C_n)\) is indexed by the divided-power
degree \(n\), so \(f\in S_n\) has redundancy \(n+1\);
\(\mathrm{SC}(j)\) is indexed directly by redundancy.  Accordingly,
after the two base cases below, \(\mathrm{SC}(j)\) is the recursive
content of
\(\mathrm{CC}(j-1,1;\mathcal P_j\cup\mathcal M_j)\).

\begin{proof}[Proof of Theorem~\ref{thm:stable-component-headline}]
We first eliminate the retained markers.  Put \(m=j-5\) and write
\[
\operatorname{Cat}_{m,4}(f)=
\begin{pmatrix}
a_0&a_1&a_2&a_3&a_4\\
a_1&a_2&a_3&a_4&a_5\\
\vdots&\vdots&\vdots&\vdots&\vdots\\
a_m&a_{m+1}&a_{m+2}&a_{m+3}&a_{m+4}
\end{pmatrix}.
\]
If \(h=(h_0,\ldots,h_m)\) is the coefficient vector of the marker
product, its bottom syndrome is
\begin{equation}\label{eq:stable-rowspace-map}
                       c(h)=h\operatorname{Cat}_{m,4}(f).
\end{equation}
Over an algebraic closure, products of \(m\) distinct linear forms are
dense in \(\PP(\Sym^mE^\vee)\).  The map in
\eqref{eq:stable-rowspace-map} is linear and homogeneous.  After deleting
its projective kernel, the scheme-theoretic closure of the attainable
bottom syndromes is therefore
\begin{equation}\label{eq:stable-rowspace}
\overline{\{c(h):h\text{ is a product of distinct markers}\}}
=\PP\!\left(\operatorname{rowspace}
              \operatorname{Cat}_{m,4}(f)\right).
\end{equation}
This includes the homogeneous boundary at infinity.  Marker diagonals
are removed only from the dense source open and reappear in its closure.
Since the right-hand side is irreducible, containment in the finite
bottom bad union implies containment in one of its irreducible
components.  We now construct and check that component ledger.

Let \(\mathbf G=\operatorname{Gr}(2,\Gamma^3E)\) over
\(\operatorname{Spec}\mathbf Z\), and write
\[
(z_0,\ldots,z_5)=(p_{01},p_{02},p_{03},p_{12},p_{13},p_{23}),
\qquad z_0z_5-z_1z_4+z_2z_3=0.
\]
After dividing the alternating evaluation determinant of a cubic pencil
by the diagonal bracket, its ordered-root incidence is
\begin{equation}\label{eq:stable-root-incidence}
\begin{aligned}
G_z(t,s)={}&z_0+z_1(t+s)+z_2(t^2+ts+s^2)+z_3ts\\
            &+z_4ts(t+s)+z_5t^2s^2 .
\end{aligned}
\end{equation}
Its homogenization is intrinsic on
\(\PP(E^\vee)\times\PP(E^\vee)\).  We retain the closed collision
incidence before taking a residual.

Over an algebraically closed field, a reducible moving part either has a
vertical or horizontal factor, hence a fixed-root or collision component,
or splits into two \((1,1)\) factors.  Symmetry under
\(\tau(t,s)=(s,t)\) gives three cases.  If both factors are symmetric,
\[
L=a+b(t+s)+cts,\qquad M=A+B(t+s)+Cts,
\]
then
\begin{equation}\label{eq:stable-symmetric-factors}
z_0z_5-z_1z_4+z_2z_3=(ac-b^2)(AC-B^2),
\end{equation}
so one factor has rank one and gives collision.  If the factors are
exchanged,
\[
L=a+bt+cs+dts,\qquad G_z=L(t,s)L(s,t),
\]
then
\begin{equation}\label{eq:stable-exchanged-factors}
z_0z_5-z_1z_4+z_2z_3
=(ad-bc)(ad-b^2+bc-c^2).
\end{equation}
The first factor is again collision; the second is the cyclic branch.
Finally, if \(\tau L=-L\), then \(L\) is proportional to the homogeneous
diagonal bracket.  The product of the two anti-invariant factors has
\[
(z_0,\ldots,z_5)=(0,0,-\mu,3\mu,0,0)
\]
and Plücker value \(-3\mu^2\).  It survives only in characteristic
three, where it belongs to the retained diagonal collision incidence.
These three cases are exhaustive by unique factorization.

For a bottom syndrome \(c=(c_0,\ldots,c_4)\), the kernel line of
\[
H_c=\begin{pmatrix}c_0&c_1&c_2&c_3\\
                    c_1&c_2&c_3&c_4\end{pmatrix}
\]
has Plücker coordinates
\begin{equation}\label{eq:stable-plucker-map}
\begin{aligned}
z_0&=c_2c_4-c_3^2,&z_1&=-(c_1c_4-c_2c_3),&
z_2&=c_1c_3-c_2^2,\\
z_3&=c_0c_4-c_1c_3,&z_4&=-(c_0c_3-c_1c_2),&
z_5&=c_0c_2-c_1^2 .
\end{aligned}
\end{equation}
Eliminating \(a,b,c,d\) from the noncollision factor of
\eqref{eq:stable-exchanged-factors} gives
\begin{equation}\label{eq:stable-J-generators}
\begin{aligned}
K_1&=3z_4^2-9z_2z_5-z_3z_5,&K_2&=z_3z_4-3z_1z_5,\\
K_3&=z_3^2-9z_0z_5,&K_4&=z_2z_3-z_1z_4+z_0z_5,\\
K_5&=z_1z_3-3z_0z_4,&K_6&=3z_1^2-9z_0z_2-z_0z_3 .
\end{aligned}
\end{equation}
Thus \(J=(K_1(z(c)),\ldots,K_6(z(c)))\).  The primitive cyclic ideal
\(V=(V_1,\ldots,V_7)\) is
\begin{equation}\label{eq:stable-V-generators}
\begin{aligned}
V_1={}&2c_3^3-3c_2c_3c_4+c_1c_4^2,\\
V_2={}&6c_2c_3^2-9c_2^2c_4+2c_1c_3c_4+c_0c_4^2,\\
V_3={}&2c_1c_3^2-3c_1c_2c_4+c_0c_3c_4,\\
V_4={}&c_0c_3^2-c_1^2c_4,\\
V_5={}&2c_1^2c_3-3c_0c_2c_3+c_0c_1c_4,\\
V_6={}&6c_1^2c_2-9c_0c_2^2+2c_0c_1c_3+c_0^2c_4,\\
V_7={}&2c_1^3-3c_0c_1c_2+c_0^2c_3 .
\end{aligned}
\end{equation}
Put
\[
D=\det\begin{pmatrix}
c_0&c_1&c_2\\c_1&c_2&c_3\\c_2&c_3&c_4
\end{pmatrix}.
\]
Direct integral expansion gives
\begin{equation}\label{eq:stable-integral-bridge}
                         J\subset V,\qquad 3DV\subset J .
\end{equation}
The supplement prints the \(6\times7\) coefficient matrix for the second
containment.  Integer reduction also gives \(DV\not\subset J\), and the
\(\mathbf F_3\)-point \((1,0,0,1,0)\) lies on \(J\), has \(D=2\), and
does not lie on \(V\).  Hence \(3\) is minimal for this elimination
presentation, not asserted to be an intrinsic invariant.

Over \(\mathbf Z[1/6]\), diagonal rescaling identifies
\eqref{eq:stable-V-generators} with the prime homogeneous kernel of
\[
[u:v:w]\longmapsto[u^2:2uv:v^2+2uw:2vw:w^2].
\]
The image point \((1,0,2,0,1)\) has \(D=-6\), so \(D\notin V\) and
\(V:D^\infty=V\).  Therefore
\begin{equation}\label{eq:stable-saturation}
(J:D^\infty)=(V:D^\infty)\ \text{over }\mathbf Z[1/3],
\qquad
(J:D^\infty)=V\ \text{over }\mathbf Z[1/6].
\end{equation}
The inversion of \(2\) in the second statement is necessary: in
characteristic two, \(V\) retains a \(D\)-supported component, so
\((V:D^\infty)\ne V\).

For a homogeneous carrier ideal \(I\), define its
\emph{coherent-Fano ideal} by substituting the universal consecutive
polar line into every generator of \(I\), expanding in the two line
parameters, and setting all coefficients to zero.  It represents upper
syndromes whose entire polar line lies in the carrier, with the required
consecutive-Hankel integrability constraint.

For \(d=(d_0,\ldots,d_5)\), substitute
\(c_i=xd_i+yd_{i+1}\) in \eqref{eq:stable-V-generators}, and let
\(F_{i,k}\) be the coefficient of \(x^{3-k}y^k\) in \(V_i(c(x,y))\).
If \(H_1,\ldots,H_4\) are the maximal minors, in column order, of the
consecutive \(3\times4\) Hankel matrix of \(d\), then
\begin{equation}\label{eq:stable-fano-certificate}
\begin{aligned}
6H_1&=-3F_{4,0}-2F_{5,1}+F_{6,2}-6F_{7,3},\\
6H_2&=-6F_{3,0}+9F_{4,1}+6F_{5,2}-3F_{6,3},\\
6H_3&=-3F_{2,0}+6F_{3,1}-9F_{4,2}-6F_{5,3},\\
6H_4&=-6F_{1,0}+F_{2,1}-2F_{3,2}+3F_{4,3}.
\end{aligned}
\end{equation}
The sign pattern is checked by reversal:
\(H_1\leftrightarrow H_4\), \(H_2\leftrightarrow H_3\),
\(V_1\leftrightarrow V_7\), \(V_2\leftrightarrow V_6\),
\(V_3\leftrightarrow V_5\), while \(V_4\mapsto-V_4\).
Thus magnitudes reverse and exactly the \(F_4\) position changes sign in
each pair.  Equation \eqref{eq:stable-fano-certificate} puts the tame
cyclic coherent-Fano locus inside the persistent locus over
\(\mathbf Z[1/6]\).

We next treat the vertical fibres.  In characteristic three the true wild
carrier is the cone \(\operatorname{Join}(e_2,C_4)\).  Any line not
through its vertex would project to a line in the quartic normal rational
curve, so every contained line is a ruling.  Move a ruling to
\(\langle e_2,e_4\rangle\).  The conditions
\[
(a_0,\ldots,a_4),(a_1,\ldots,a_5)\in\langle e_2,e_4\rangle
\]
give \(a_0=a_1=a_3=0\) and \(a_1=a_2=a_4=0\), leaving only \(a_5\).
Undoing the frame gives \(f=\mu\nu_5(t)\), including infinity, hence
the rank-one/fixed-factor boundary.

In characteristic two the true cyclic component is
\(\Pi_{\rm cyc}=(c_0,c_4)\); the other component is supported on \(D=0\).
The remaining inseparable degeneration is seen on the
\emph{ordered-Hessian chart}: for an ordered pencil basis \(p_0,p_1\),
it is the bidegree-\((2,2)\) equation
\[
u^2\operatorname{Hess}^{\rm div}(p_0)
+uv\,\operatorname{PolHess}^{\rm div}(p_0,p_1)
+v^2\operatorname{Hess}^{\rm div}(p_1)=0.
\]
The tangent quadric has two conic families of generator lines, called
the tangent-quadric rulings:
\[
\begin{aligned}
z_1=z_4=0,\quad z_2=z_3,\quad z_0z_5=z_2^2,\\
z_0=z_5=0,\quad z_2=z_3,\quad z_1z_4=z_2^2.
\end{aligned}
\]
The first is the square-quadratic boundary of the persistent Veronese.
For the complementary ruling, the two consecutive identities
\[
h_{-1}^2=h_{-2}h_0,\qquad h_1^2=h_0h_2
\]
propagate across their common term and make all five contraction forms
proportional, hence rank one.  Vertical and horizontal factors are,
respectively, collision and fixed-factor components.

It remains to apply the rowspace reduction
\eqref{eq:stable-rowspace}.  If the row space lies in the rank-two
Hankel component, induct on \(m\): every first contraction has its
\((m-1)\)-marker row space in the same component, and
Proposition~\ref{prop:contained-rank-two} then puts \(f\) in the upper
persistent scheme.  For a linear modular component
\(L_S=\langle e_i:i\in S\rangle\subset\Gamma^dE\), the two contractions
of an upper basis vector \(e_q\) are, up to divided-power normalization,
\(e_q\) and \(e_{q-1}\).  Consequently its coherent lift is the
scheme-theoretic kernel
\[
\ker\!\left(
 \Gamma^{d+1}E\longrightarrow
 \operatorname{Hom}(E^\vee,\Gamma^dE/L_S)
\right)
=\langle e_q:q,q-1\in S\rangle .
\]
After \(\ell\) lifts its support is therefore exactly the set of \(q\)
for which \(q,q-1,\ldots,q-\ell\in S\).  For an NRC nucleus, Lucas'
theorem computes \(S\) from the vanishing binomial digits.  Thus the
coherent lifts are precisely the Lucas consecutive-support components:
the kernel description is scheme-theoretic, commutes with base change,
and leaves neither an extra nonlinear component nor an embedded one.
The tame projected Veronese contains no line, so its
rowspace pullback has rank one.  A row space in the characteristic-three
cone is a ruling or a point; applying the displayed two-row overlap to
each adjacent pair again gives rank one.

Finally, in characteristic two, rowspace containment in
\(\Pi_{\rm cyc}\) forces the first and last columns of the catalecticant
to vanish:
\[
a_0=\cdots=a_m=0,\qquad a_4=\cdots=a_{m+4}=0.
\]
For \(m=1,2\) these are the known \(N_3\) and \([e_3]\) pullbacks.
For \(m\geq3\) the blocks cover every coefficient, so no projective
point remains.  Hence the characteristic-two modular component descended
from the R5 cyclic plane terminates after redundancy seven.  This does
not erase the fresh higher Lucas carriers of
the separate Lucas-carrier analysis; for example
\(\PP\langle e_2,\ldots,e_7\rangle\subset\Gamma^9E\) is nonempty.

The collision and fixed-factor alternatives are already covered by
Lemma~\ref{lem:uniform-collision} and the gcd/rank boundary.
This exhausts the finite bottom component ledger and proves
\(\mathrm{SC}(j)\) for every \(j\geq6\), together with the asserted
termination of the cyclic-plane descendant.
\end{proof}

The rank-one boundary, ordinary rank-two component, collision
alternative, and modular pullback are uniform by
Proposition~\ref{prop:contained-rank-two} and
Lemmas~\ref{lem:uniform-collision} and
\ref{lem:linear-modular-pullback}.  The two inseparable collision
possibilities are exactly the characteristic-three R5 case and the
characteristic-two R6 case treated in
Propositions~\ref{prop:r5-bridge} and \ref{prop:r6-degrees}.
Section~\ref{sec:r67} independently recovers the redundancies-six and
seven base cases.

\begin{theorem}[One-step polar escape]\label{thm:induction}
Let \(f\in S_n(\F_q)\) have a one-step lower package as in
Definition~\ref{def:lower-package}.  Put
\[
\mathcal H_\kappa(g,\delta)=
1+\left\lfloor
\left(g+\sqrt{g^2+\delta-\kappa}\right)^2
\right\rfloor
\]
and assume
\(\delta_\alpha\geq\kappa_\alpha-g_\alpha^2\) for every \(\alpha\).
If
\[
q\geq \max_\alpha
\mathcal H_{\kappa_\alpha}(g_\alpha,\delta_\alpha)
\quad\text{and}\quad q+1>d,
\]
then \(W_f\) contains a split squarefree member.
\end{theorem}

\begin{proof}
Since \(\#\PP^1(\F_q)=q+1>d\), choose
\(\lambda\notin Z_f(\F_q)\).  On the corresponding lower stratum the
recorded point bound
gives
\[
 q+\kappa_\alpha-2g_\alpha\sqrt q>\delta_\alpha,
\]
so there is a rational point outside all deletions.  It yields a split
squarefree lower member avoiding \(\lambda\).  Equation
\eqref{eq:lift} lifts it to a split squarefree member upstairs,
so the original system is shallow.
\end{proof}

The theorem isolates the numerical escape and one-step lifting
argument.  Iterating it requires the contained-component assertion at
every intermediate level, not only at the top.

Recursively, the current contraction is chosen outside
\(\mathcal P_j\cup\mathcal M_j\); the uniform alternatives and the next
\(\mathrm{SC}\) theorem make its polar line transverse, so a new
marker exists off the finite pullback.  No inverse-image stability of
the declared loci is assumed.

\begin{remark}[Point-count convention]\label{rem:thresholds}
The correction $\kappa$ records the exact lower bound available for the
chosen model.  Every genus-one curve retained in this paper uses the
Aubry--Perret bound \cite[p.~468]{AubryPerret1995} with \(\kappa=1\);
a possible rational singular
point is included instead in the deletion degree.
\end{remark}

\begin{theorem}[Uniform transverse threshold]
\label{thm:uniform-transverse}
Fix a redundancy \(r\geq6\).  If
\[
 q\geq Q_r:=
 6r-15+\left\lfloor2\sqrt{6r-17}\right\rfloor ,
\]
every split-free redundancy-\(r\) syndrome belongs to
\(\mathcal P_r\cup\mathcal M_r\).
\end{theorem}

\begin{proof}
Iterate the pointed contraction to the R5 cubic pencil.  Its
off-diagonal curve always has \((g,\kappa)=(1,1)\), and its singular,
diagonal, and branch deletions cost thirteen.  Each of the \(r-5\)
retained markers removes at most six further ordered pairs.  Hence
\[
 \delta_r=13+6(r-5)=6r-17,\qquad
 \mathcal H_1(1,\delta_r)
 =6r-15+\left\lfloor2\sqrt{6r-17}\right\rfloor .
\]
At a stage of redundancy \(j\) with \(s=r-j\) old markers, the
ordinary pullback costs at most three and each old marker at most
three.  At the bottom stage \(j=6\), the R5 cyclic/wild pullback costs
four by Propositions~\ref{prop:r6-degrees} and
\ref{prop:r6-modular}, while the collision divisor costs six, giving
\[
 d_6\leq3+4+6+3(r-6)=3r-5.
\]
For \(j\geq7\), the recursive bad scheme and \(\mathrm{SC}(j)\) leave
only the persistent and modular lower components.
Lemmas~\ref{lem:uniform-collision} and
\ref{lem:linear-modular-pullback} give
\[
 d_j\leq3+1+(2j-6)+3(r-j)=3r-j-2<3r-5.
\]
Thus the maximum is attained at the bottom stage, and \(Q_r>3r-5\)
makes the parameter-line inequality \(q+1>d_j\) automatic.  The
exact-linear-gcd graph has eight base deletions and at most two per
retained marker, hence
\[
 \delta_r^{\mathrm{lin}}=8+2(r-5)=2r-2,
\qquad
 Q_r+1-\delta_r^{\mathrm{lin}}
 =4r-12+\left\lfloor2\sqrt{6r-17}\right\rfloor>0.
\]
Thus this comparison is weaker for every \(r\geq6\).

Apply Theorem~\ref{thm:induction} successively.  If a recursively
pointed bad scheme contains the whole polar line, the uniform
alternatives above and \(\mathrm{SC}(j)\) place the syndrome in the
declared persistent or modular list.  Otherwise a marker remains
outside its finite scheme.  After reaching the R5 pencil, the
displayed point bound supplies a split squarefree cubic avoiding every
marker, and Theorem~\ref{thm:polar-construction} lifts it to the
original system.
\end{proof}

For \(r=5,6,7\), the integer cutoffs supplied by the same formula are
\[
 Q_5=22,\qquad Q_6=29,\qquad Q_7=37;
\]
the first prime powers meeting them are \(23,29,37\).  Thus the former
threshold table consisted of three evaluations of one formula with
\(Q_r=6r+O(\sqrt r)\).  The same calculation recovers the verified
R6 and R7 parameter bounds \(13\) and \(16\).  In general the
parameter line is nonbinding; the stable-component assertion is supplied
uniformly by Theorem~\ref{thm:stable-component-headline}.
Propositions~\ref{prop:r6-contained} and
Corollary~\ref{cor:r7-contained} discharge the required assertions at
redundancies six and seven.

\subsection{Contained-component calculations}

The ordinary persistent locus admits a uniform contained-line
calculation.  This does not classify other lower bad schemes, but it
closes the persistent part of every level-specific component gate.

\begin{proposition}[Contained rank-two polar lines]
\label{prop:contained-rank-two}
Let $n\geq5$, let $f\in\PP(\Gamma^nE)$, and let
\[
 A=C_f^{(2)}:\Sym^{n-2}E^\vee\longrightarrow\Gamma^2E
\]
be the three-row catalecticant.  Assume the contraction map
$E^\vee\to\Gamma^{n-1}E$ is injective, so the first-polar map is
defined on all of $\PP(E^\vee)$.  This is exactly the complement of
the rank-one boundary: after a change of coordinates, a nonzero
kernel vector is $\partial_X$, and
$\partial_Xf=0$ forces $f$ to be a scalar multiple of $Y^{[n]}$.
Such a point lies on the normal rational curve and is already excluded
from the split-free locus.  Then
\[
\Phi_f(\PP(E^\vee))\subset
\{b:\operatorname{rank}C_b^{(2)}\leq2\}
\quad\Longleftrightarrow\quad
\operatorname{rank}A\leq2
\]
scheme-theoretically and in every characteristic.  If
$\operatorname{rank}A=3$, the rank-drop parameters form a subscheme of
$\PP(E^\vee)$ of degree at most three.
\end{proposition}

\begin{proof}
For a nonzero $\lambda\in E^\vee$, multiplication identifies
$\Sym^{n-3}E^\vee$ with the hyperplane
\[
 H_\lambda=\lambda\Sym^{n-3}E^\vee\subset\Sym^{n-2}E^\vee,
\]
and functoriality gives
\[
 C_{\iota_\lambda f}^{(2)}=A|_{H_\lambda}.
\]
The forward implication is immediate.  Conversely, suppose
$\operatorname{rank}A=3$ and put $K=\ker A$, so
$\dim K=n-4$.  If the lower rank is at most two, rank--nullity on the
$(n-2)$-dimensional space $H_\lambda$ gives
\[
 \dim(K\cap H_\lambda)\geq(n-2)-2=n-4=\dim K.
\]
Thus $K\subset H_\lambda$.  Containment for every geometric
$[\lambda]$ would imply
\[
 K\subset\bigcap_{[\lambda]\in\PP(E^\vee_{\bar k})}
 \lambda\Sym^{n-3}E^\vee_{\bar k}=0,
\]
because a nonzero binary form cannot be divisible by every linear
form.  This contradicts $\dim K=n-4>0$.

The maximal minors of $C_{\iota_\lambda f}^{(2)}$ are homogeneous
cubics in $\lambda$.  Vanishing at every geometric point is therefore
identical vanishing, proving the scheme-theoretic assertion.  In the
noncontained case at least one cubic is nonzero, and the common
rank-drop scheme lies in its degree-three zero divisor.
\end{proof}

Thus Proposition~\ref{prop:contained-rank-two} covers the entire
non-rank-boundary locus on which the first-polar line is defined; the
excluded noninjective locus is the rank-one normal-rational-curve
locus just identified.  After rank one and rational split secants are removed as shallow, the
upper rank-two locus consists of the persistent tangent and
conjugate-secant families.  On a persistent fiber the nonsplit torus
acts by $z\mapsto z^n$, giving $T/T^n$ modulo inversion and Frobenius,
while the tangent unipotent cocycle is $z\mapsto z+nu$.

Modular loci are equally explicit.  If $\cN$ is a nucleus in the
lower divided-power module, the complete space whose first-polar line
lies in $\cN$ is
\[
\cM_n(\cN)=
\PP\ker\left(
\Gamma^nE\longrightarrow
\operatorname{Hom}(E^\vee,\Gamma^{n-1}E/\widetilde{\cN})
\right).
\]
Lucas' theorem computes the nucleus coordinates, and this displayed
kernel computes every coherent lift without an ambient syndrome census.
At redundancy \(r\), if \(p>r-1\), Lucas' theorem leaves no nucleus
coordinate, so \(\mathcal M_r\) is empty.  This removes the modular
term only; the R5 cyclic carrier survives in large characteristic.

\begin{remark}[Scope of the contained calculations]
The modular contraction kernel is an exact finite linear-algebra
construction.  Proposition~\ref{prop:contained-rank-two} supplies the
ordinary persistent component uniformly, and
Lemma~\ref{lem:uniform-collision} supplies the collision alternative.
The remaining arbitrary-level question is \(\mathrm{SC}(j)\): whether
another contained component occurs.
\end{remark}
